\section*{अध्याय \ref{ch:b3:glaciology} --- हिमविज्ञान}

\begin{solution}{exo:b3:glaciology:1}
$\tau_b = \rho gH\sin\alpha = 917\times9.81\times250\times\sin\ang{3} =
\qty{118}{kPa}$, लगभग \qty{1.2}{bar}, यानी वास्तविक हिमनदों के नीचे
के प्रतिबलों की कोटि।
\end{solution}

\begin{solution}{exo:b3:glaciology:2}
(क)~$\qty{2.9e15}{m^3}\times917 = \qty{2.66e18}{kg} = \qty{2.66e6}{Gt}$;
\qty{360}{Gt/mm} से भाग देने पर \qty{7390}{mm}, अर्थात \qty{7.4}{m}, जो उद्धृत मान है।
(ख)~$\qty{3.0e16}{m^3}\times917/360 = \qty{76}{m}$, जबकि उद्धृत मान \qty{58}{m} है।
ग्रीनलैंड की बर्फ़ अधिकतर समुद्र-तल से ऊपर के आधार-शैल पर टिकी है, इसलिए उसका सब कुछ
पिघलकर महासागर में पानी जोड़ेगा। अंटार्कटिका का अधिकांश हिस्सा समुद्र-तल से नीचे के तल पर
टिका है: वह बर्फ़ पहले ही समुद्री जल को विस्थापित करती है, इसलिए उसके पिघलने से
केवल वही जल भरता है जिसे वह विस्थापित करती है; तैरते हिम-शेल्फ़ तो बिल्कुल कुछ नहीं जोड़ते।
\end{solution}

\begin{solution}{exo:b3:glaciology:3}
लाभ $\num{6e6}\times1.2 = \qty{7.2e6}{m^3}$, हानि $\num{4e6}\times2.1 =
\qty{8.4e6}{m^3}$ जल: शुद्ध \qty{-1.2e6}{m^3}, माध्य विशिष्ट संतुलन
$-1.2\times10^6/10^7 = \qty{-0.12}{m}$ जल-तुल्यांक होता है।
संचयन-क्षेत्र अनुपात $6/10 = 0.6$ है। हिमनद द्रव्यमान खो रहा है:
सिकुड़ रहा है।
\end{solution}

\begin{solution}{exo:b3:glaciology:4}
$1 - 820/917 = \num{0.106}$ और $1 - 830/917 = \num{0.095}$:
आयतन का लगभग एक-दसवाँ भाग अब भी रंध्र-स्थान है, जब रंध्र बंद होते हैं।
\end{solution}

\begin{solution}{exo:b3:glaciology:5}
$\tau_b = 917\times9.81\times300\times\sin\ang{3} = \qty{141.2}{kPa}$। स्लैब प्रमेय से,
समीकरण प्रमेय से, $u_s = \tfrac{2A}{4}\tau_b^3H = 0.5\times\num{1.43e-16}
\times(\num{1.412e5})^3\times300 = \qty{60}{m/yr}$। दोगुनी मोटाई होने पर $u_s
\propto H^{n+1} = H^4$, सोलह गुना तेज़। ढाल आधी होने पर $u_s \propto
(\sin\alpha)^3$, आठ गुना धीमी।
\end{solution}

\begin{solution}{exo:b3:glaciology:6}
\qty{-25}{\celsius} (\qty{248.15}{K}) पर, $A = \num{3.61e-13}\exp(-60000/
(8.314\times248.15)) = \qty{8.5e-26}{Pa^{-3}.s^{-1}}$, या
\qty{2.7e-18}{Pa^{-3}.yr^{-1}}। \qty{-10}{\celsius} पर, $\num{3.61e-13}
\exp(-27.42) = \qty{4.4e-25}{Pa^{-3}.s^{-1}} = \qty{1.4e-17}{Pa^{-3}.yr^{-1}}$।
\qty{0}{\celsius} पर (गर्म शाखा) \qty{1.4e-16}{Pa^{-3}.yr^{-1}}। समान प्रतिबलों के अंतर्गत
गतियाँ दर गुणकों के अनुपात में होती हैं:
$\num{1.43e-16}/\num{2.67e-18} = 54$।
\end{solution}

\begin{solution}{exo:b3:glaciology:7}
$h = \sqrt{2\times10^5\times10^6/(917\times9.81)} = \qty{4.7}{km}$, जो अंटार्कटिका की
\qty{4.9}{km} मोटाई के निकट है: प्लास्टिक मॉडल किसी हिमचादर की मोटाई
उसकी चौड़ाई से ही प्राप्त कर लेता है।
\end{solution}

\begin{solution}{exo:b3:glaciology:8}
$\tau = 200/5 = \num{40}\text{ वर्ष}$। ऊष्मीकरण रुकने के बाद लंबाई
कई दशकों, यानी कुछ प्रतिक्रिया-काल तक, समायोजित होती रहेगी: आज के हिमनद
अभी तक आज की जलवायु से तालमेल नहीं बिठा पाए हैं।
\end{solution}

\begin{solution}{exo:b3:glaciology:9}
(क)~$h = \sqrt{2\times2.03\times20\times\num{8.64e6}/(917\times\num{3.34e5})}
= \qty{1.51}{m}$। चूँकि $h^2 \propto t$, \qty{3}{m} की मोटाई तक पहुँचने में $100\times
(3/1.51)^2 = 393$ दिन लगते हैं, जो एक सर्दी से अधिक है। (ख)~
\qty{3}{m} से अधिक मोटी बर्फ़ एक सर्दी में जमकर नहीं बढ़ सकती: वह ऐसी बर्फ़ है जो
कई सर्दियों तक बची रही है, या हवा और धाराओं द्वारा कटकों में जमा की गई बर्फ़।
\end{solution}

\begin{solution}{exo:b3:glaciology:10}
(क)~$H = (1025/917)\times800 = \qty{894}{m}$। (ख)~\qty{900}{m} पर
\omterm{def:b3:glaciology:grounding-line}{भू-स्पर्श रेखा} की मोटाई \qty{1006}{m} है, पूर्व मान का $9/8$ गुना; फ़्लक्स
$(9/8)^5 = 1.8$ गुना बढ़ जाता है। (ग)~जैसे-जैसे रेखा भीतर की ओर गहराते तल पर पीछे हटती है,
रेखा पर की बर्फ़ अधिक मोटी होती है, अधिक विसर्जन करती है, और पीछे की
बर्फ़ को पतला करती है, और रेखा और पीछे हटती है: यह धनात्मक प्रतिपुष्टि है। हिम-शेल्फ़
तल पर टिकी बर्फ़ को सहारा देते हुए \omterm{def:b3:glaciology:grounding-line}{भू-स्पर्श रेखा} पर फ़्लक्स घटाता है और
उसे नीचे की ओर ढलान पर भी उसी स्थान पर रोक सकता है।
\end{solution}

\begin{solution}{exo:b3:glaciology:11}
जहाँ $h(r) = c\sqrt{R - r}$, $c = \sqrt{2\tau_0/(\rho g)}$, $V =
\int_0^R2\pi rc\sqrt{R - r}\,\mathrm dr$। $s = R - r$ प्रतिस्थापित करने पर,
\[
\int_0^R(R - s)\sqrt s\,\mathrm ds = \tfrac23R^{5/2} - \tfrac25R^{5/2} =
\tfrac4{15}R^{5/2},
\]
अतः $V = \frac{8\pi}{15}cR^{5/2}$। $R$ को दोगुना करने पर
$V$ और \omterm{def:b3:glaciology:sle}{समुद्र-तल तुल्यांक}, $2^{5/2} = 5.7$ गुना हो जाते हैं।
\end{solution}

\begin{solution}{exo:b3:glaciology:12}
(क)~मार्च: $14.13/16.34 - 1 = \qty{-13.5}{\percent}$; सितंबर: $4.81/7.05
- 1 = \qty{-31.8}{\percent}$। (ख)~मार्च में हिम-किनारा अधिकतर उन तटों से थमा रहता है जो
आर्कटिक महासागर को घेरते हैं, और हर सर्दी का पुनर्जमाव
बेसिन को फिर से भर देता है; हाल के वर्षों की पतली बर्फ़ हर गर्मी में अधिक पूरी तरह
पिघल जाती है। और हिम--एल्बिडो प्रतिपुष्टि गर्मी में काम करती है: खुला पानी
लगभग \qty{94}{\percent} सौर प्रकाश सोखता है, बर्फ़ परावर्तित करती है
\qtyrange{50}{70}{\percent} के बीच, इसलिए गर्मी की हर बर्फ़-हानि जल को गर्म करती है और
और अधिक पिघलती है।
\end{solution}

\begin{solution}{pb:b3:glaciology:1}
\textbf{1.} तल के ऊपर का स्तंभ, जिसका प्रति इकाई क्षेत्र पर भार $\rho gH$ है,
ढाल की दिशा में नीचे की ओर घटक $\rho gH\sin\alpha$ से धकेला जाता है, जिसे
उसके आधार पर का अपरूपण प्रतिबल संतुलित करता है; छोटे ढलानों के लिए $\sin\alpha \approx \alpha$ (
पृष्ठ ढलान के साथ स्लैब प्रमेय)।

\textbf{2.} $917\times9.81\times580\times0.0190 = \qty{99.1}{kPa}$;
$917\times9.81\times470\times0.0240 = \qty{101.5}{kPa}$; $917\times9.81
\times335\times0.0330 = \qty{99.4}{kPa}$।

\textbf{3.} तनाव \qty{2}{\percent} के भीतर समान रहता है, जबकि
मोटाई लगभग आधी हो जाती है: टोपी लगभग सुघट्य है, जिसका $\tau_0 \approx
\qty{100}{kPa}$ है।

\textbf{4.} $2\times10^5/(917\times9.81) = \qty{22.23}{m}$।

\textbf{5.} $h = \sqrt{22.23\,(R - r)}$: किनारे से \qty{15}{km}, \qty{10}{km} और \qty{5}{km} पर क्रमशः \qty{577}{m}, \qty{472}{m} और
\qty{333}{m} मोटाई मिलती है, जो
सर्वेक्षण-मापों से कुछ मीटर के भीतर है।

\textbf{6.} $\sqrt{22.23\times20000} = \qty{667}{m}$।

\textbf{7.} \cref{exo:b3:glaciology:11} के समान: $V = \int_0^R2\pi r
\sqrt{22.23(R - r)}\,\mathrm dr = \frac{8\pi}{15}\sqrt{2\tau_0/(\rho
g)}\,R^{5/2}$।

\textbf{8.} $\frac{8\pi}{15}\times4.715\times(\num{2e4})^{5/2} =
\qty{4.47e11}{m^3} = \qty{447}{km^3}$।

\textbf{9.} $\pi\times20^2 = \qty{1257}{km^2}$; माध्य मोटाई
$447/1257 = \qty{0.356}{km}$, यानी \qty{356}{m}, जो
केंद्रीय मोटाई के आधे से थोड़ा अधिक है।

\textbf{10.} $2^{5/2} = 5.7$ गुना।

\textbf{11.} $A = \num{3.61e-13}\exp(-60000/(8.314\times263.15)) =
\qty{4.45e-25}{Pa^{-3}.s^{-1}}$, को \qty{3.156e7}{s} से गुणा करने पर:
\qty{1.40e-17}{Pa^{-3}.yr^{-1}}।

\textbf{12.} \qty{10}{km} पर, $H = \qty{471.5}{m}$, ढाल $471.5/20000 =
\num{0.0236}$, $\tau_b = \qty{100}{kPa}$; $u_s = \tfrac{2A}{4}\tau_b^3H =
0.5\times\num{1.40e-17}\times10^{15}\times471.5 = \qty{3.3}{m/yr}$।

\textbf{13.} $q = \tfrac{2A}{5}\tau_b^3H^2 = 0.4\times\num{1.40e-17}
\times10^{15}\times471.5^2 = \qty{1250}{m^2/yr}$; माध्य वेग $q/H =
\qty{2.6}{m/yr}$, जो पृष्ठ की गति का चार-पाँचवाँ भाग है।

\textbf{14.} $2\pi rq = \pi r^2\bar b$, अतः $\bar b = 2q/r = 2\times1250/10^4
= \qty{0.25}{m}$ हिम प्रति वर्ष, भीतरी \qty{10}{km} पर।

\textbf{15.} $h = 450$ उस स्थान पर, जहाँ $R - r = 450^2/22.23 = \qty{9.11}{km}$ है, अर्थात
$r = \qty{10.9}{km}$ पर; \omterm{def:b3:glaciology:mass-balance}{संचयन क्षेत्र} इसके भीतर की डिस्क है, $(10.9/
20)^2 = 0.30$ भाग टोपी का।

\textbf{16.} $A$ \num{10.2} गुना बढ़ता है, और उसी प्रकार समान प्रतिबलों के अंतर्गत गतियाँ भी।
हिम, संतुलन जितनी आपूर्ति करता है उससे तेज़ बाहर बहेगा: टोपी पतली होगी और फैलेगी,
जब तक घटे हुए प्रतिबल फिर से संतुलन से मेल न खा जाएँ;
एक नरम सुघट्य टोपी का विराम प्रतिबल छोटा होता है और
उसकी प्रोफ़ाइल अधिक समतल होती है।

\textbf{17.} $\num{4.47e11}\times917 = \qty{4.10e14}{kg} = \qty{410}{Gt}$।

\textbf{18.} $410/360 = \qty{1.14}{mm}$।

\textbf{19.} $356/0.5 = 711$ वर्ष।

\textbf{20.} जैसे-जैसे सतह नीचे होती है, वह गर्म वायु में पड़ती है और तेज़ी से पिघलती है,
ऊँचाई--द्रव्यमान संतुलन प्रतिपुष्टि के कारण; और जैसे-जैसे वह पतली होती है, प्रतिबल तथा
प्रवाह भी बदलते हैं। पतलापन तेज़ होता जाता है।

\textbf{21.} $410/243 = \num{1.7}\text{ वर्ष}$: पूरी टोपी 2010--2019 में ग्रीनलैंड
की हानि के दो वर्षों से कम है।

\textbf{22.} आर्केल हिमटोपी का \omterm{def:b3:glaciology:sle}{समुद्र-तल तुल्यांक}
\textbf{\qty[text-series-to-math=true]{1.1}{mm}} (\qty{1.14}{mm}) है।
\end{solution}
